\section{Duality}\label{sec:duality}

Let $\mathsf E$ be a metrized complex of representations of $Q$ with charges
$z_i$. The \emph{dual complex} $\mathsf E^\vee$ is the complex of
representations of the opposite quiver $Q^{\mathrm{op}}$ with
\[
 (E^\vee)_i^n=(E_i^{-n})^*,\qquad d_i^\vee=d_i^{\mathsf T},\qquad
 \rho^\vee(a^{\mathrm{op}})=\rho(a)^{\mathsf T},
\]
the dual metrics, and the charges $z_i^\vee=-\overline{z_i}$, which again lie
in the upper half-plane. Here $X^{\mathsf T}$ is the transposed map between
the dual spaces; in orthonormal frames and their dual frames it is the
transposed matrix, and $\overline X$ denotes the entrywise complex conjugate.

\begin{theorem}[Duality]\label{thm:duality}
For every metrized complex $\mathsf E$:
\begin{enumerate}
\item[(a)] As endomorphisms of $E^*$,
 $\Omega^\vee=-(\Omega^\dagger)^{\mathsf T}$, $|\Omega^\vee|=|\Omega|^{\mathsf T}$ and
 $\Phi^\vee=\pi-\Phi^{\mathsf T}$; in dual orthonormal frames,
 $\Omega^\vee=-\overline\Omega$ and $\Phi^\vee=\pi-\overline\Phi$. Hence
 $\mass(\mathsf E^\vee)=\mass(\mathsf E)$ and
 $Z(\mathsf E^\vee)=-\overline{Z(\mathsf E)}$.
\item[(b)] If $h(t)$ is a trajectory of the phase flow of $\mathsf E$, then
 $e^{-\pi t}(h(t)^{-1})^{\mathsf T}$ is a trajectory of the phase flow of
 $\mathsf E^\vee$. In orthonormal frames, transposition of all structure maps
 intertwines the equations \eqref{eq:frameflow} of the two complexes.
\item[(c)] If $Z(\mathsf E)\neq0$ and $\varphi^\vee=\pi-\varphi$, the phase
 potentials $\Ima(e^{-\ii\varphi^\vee}\cS^\vee_{\CC})$ and
 $\Ima(e^{-\ii\varphi}\cS_{\CC})$ agree, up to an additive constant, on
 corresponding metrics $h$ and $c\,(h^{-1})^{\mathsf T}$, $c>0$, and the
 contributions $\cD_i^n$ of corresponding blocks agree.
\end{enumerate}
\end{theorem}

\begin{proof}
(a) Use Proposition~\ref{prop:coefficient} for $\mathsf E^\vee$ in dual
orthonormal frames, where $(d^\vee)^\dagger=\overline d$ and
$(\rho^\vee)^\dagger=\overline\rho$. Then
$\Delta_i^\vee=d_i^{\mathsf T}\overline{d_i}+\overline{d_i}d_i^{\mathsf T}=\Delta_i^{\mathsf T}$. The arrows
of $Q^{\mathrm{op}}$ out of $i$ are the arrows of $Q$ into $i$, so
\[
 \Gamma_i^\vee=-\overline{z_i}+\Bigl(\sum_{a:k\to i}\rho(a)\rho(a)^\dagger
 -\sum_{a:i\to j}\rho(a)^\dagger\rho(a)\Bigr)^{\mathsf T}=-\overline{\Gamma_i},
\]
the bracket being Hermitian. For $a:i\to j$ the arrow $a^{\mathrm{op}}$ goes
from $j$ to $i$, and
$M^\vee_{a^{\mathrm{op}}}=\overline{d_i}\,\rho(a)^{\mathsf T}-\rho(a)^{\mathsf T}\overline{d_j}=-M_a^{\mathsf T}$,
$(M^\vee_{a^{\mathrm{op}}})^\dagger=-\overline{M_a}$. In
Proposition~\ref{prop:coefficient} for $Q^{\mathrm{op}}$ the term of
$a^{\mathrm{op}}$ at the vertex $i$ is of incoming type,
\[
 +\mathcal J_{\Delta_i^{\mathsf T},\Delta_j^{\mathsf T}}\bigl(M_a^{\mathsf T},\overline{M_a}\bigr)
 =\bigl(\mathcal J_{\Delta_i,\Delta_j}(M_a^\dagger,M_a)\bigr)^{\mathsf T}
 =\overline{\mathcal J_{\Delta_i,\Delta_j}(M_a^\dagger,M_a)},
\]
which enters $\mathsf E$ at $i$ with the sign $-$; likewise the term at $j$
is $-\overline{\mathcal J_{\Delta_j,\Delta_i}(M_a,M_a^\dagger)}$, which
enters $\mathsf E$ with $+$. Finally
$K_{\Delta^{\mathsf T}}(-\overline\Gamma)=-\bigl(K_\Delta(\Gamma^\dagger)\bigr)^{\mathsf T}
=-\overline{K_\Delta(\Gamma)}$. Altogether
$\gr\,\Omega^\vee=-\overline{\gr\,\Omega}$, and the degrees $n$
and $-n$ have the same parity. If $\Omega=|\Omega|e^{\ii\Phi}$ on a block of
degree $n$, then
$-\overline\Omega=\overline{|\Omega|}\,e^{\ii(\pi-\overline\Phi)}$ with
$\overline{|\Omega|}>0$, and the spectrum of $\pi-\overline\Phi$ lies in
$(\pi n,\pi(1+n))$, the branch of degree $-n$.

(b) Transposing \eqref{eq:frameflow} and using $\Phi^{\mathsf T}=\overline\Phi$ gives
$\frac d{dt}T^{\mathsf T}=\frac12(\overline{\Phi_{\mathrm{source}}}\,T^{\mathsf T}-T^{\mathsf T}\,\overline{\Phi_{\mathrm{target}}})$.
In $\mathsf E^\vee$ the map $T^{\mathsf T}$ goes from the dual of the target to the dual
of the source, with phases $\pi-\overline{\Phi_{\mathrm{target}}}$ and
$\pi-\overline{\Phi_{\mathrm{source}}}$, and its equation
\eqref{eq:frameflow} is the same expression. For metrics: if
$g=c\,(h^{-1})^{\mathsf T}$, then $g^{-1}\dot g=\dot c/c+\Phi^{\mathsf T}$, and
$-\Phi^\vee=\Phi^{\mathsf T}-\pi$ gives $c=e^{-\pi t}$ up to a constant.

(c) For corresponding variations $\xi^\vee=-\overline\xi$,
$\tr(\Omega^\vee\xi^\vee)=\overline{\tr(\Omega\xi)}$, so
$\cS^\vee_{\CC}=\overline{\cS_{\CC}}$ up to a constant on the metrics
$h\leftrightarrow(h^{-1})^{\mathsf T}$, and
$\Ima(e^{-\ii(\pi-\varphi)}\overline{\cS_{\CC}})=\Ima(e^{-\ii\varphi}\cS_{\CC})$.
For a common factor $c>0$ the endomorphism $\Omega^\vee$ is unchanged and
$\cS^\vee_{\CC}(cg)=\cS^\vee_{\CC}(g)+Z(\mathsf E^\vee)\log c$; since
$e^{-\ii\varphi^\vee}Z(\mathsf E^\vee)$ is real, the phase potential is
invariant under common rescaling.
\end{proof}

\begin{remark}[Injective differentials]\label{rem:injective}
Duality exchanges complexes that are $V\xrightarrow{D}U$ in degrees $0,1$ at
a vertex $i_0$, with $D$ surjective, and a vector space in degree zero at the
other vertices, with complexes that are $U'\xrightarrow{D'}V'$ in degrees
$-1,0$ at $i_0$, with $D'$ injective. The cohomology representation of the
dual is the dual of the cohomology representation, with the dual metrics
(the quotient metric on $V'/D'U'$ is the metric of the orthogonal complement
of $D'U'$); subrepresentations correspond to quotients and arguments $\psi$
of central charges to $\pi-\psi$, so semistability, stability and
polystability correspond. By Theorem~\ref{thm:duality}, every statement
about the flow of the first kind of complex therefore gives one for the
second kind, with the charges $-\overline{z_i}$, the argument $\pi-\varphi$
of the central charge, and the metrics $e^{-\pi t}(h^{-1})^{\mathsf T}$; in
particular the rescaled metrics $e^{\varphi t}h(t)$ of the complex
correspond to the rescaled metrics $e^{(\pi-\varphi)t}h^\vee(t)$ of the dual,
so that the flow converges on the cohomology for one exactly when it does
for the other, and the acyclic part decays for one exactly when it does for
the other. This is how Theorem~\ref{thm:onevertex}(c) follows from
Theorem~\ref{thm:onevertex}(b), and how the decay and convergence theorems
of \cite{wue26b,wue26c} pass to complexes with an injective differential.

For $Q=K_m$ one has $Q^{\mathrm{op}}\cong Q$ by exchanging the two vertices,
and the dual charges are $(-\overline{z_1},-\overline{z_0})$. The pair
$(-1+\ii,1+\ii)$ is self-dual, and strictly ordered charges are mapped to
strictly ordered charges. For these self-dual charges the sufficient
conditions of Theorem~\ref{thm:onevertex}(b) and (c) cover different
dimension vectors when the acyclic part has rank at least two: for
cohomology of dimension vector $(k,l)$, part (b) applies to a surjective
differential at the vertex $0$ when
$\arg((l-k)+\ii(k+l))\geq\frac{\pi}3$, and part (c) to an injective
differential at the vertex $1$ when this argument is at most
$\frac{2\pi}3$; for $k=l$ both apply.
\end{remark}
